\recipeTitle{Plumcake Superbuonissimo}

\recipeSubtitle{Un plumcake super soffice, semplice e perfetto per la colazione.}

\recipeAuthor{Cinzia Bettio}

\recipeStory{Ho comprato uno stampo da plumcake, e l'ho fatto.}

% Add a real image in recipes/plumcake-superbuonissimo/images/ and uncomment the line below.
% \recipePhoto{recipes/plumcake-superbuonissimo/images/plumcake-superbuonissimo.jpg}

\recipeInfo
  {15 minuti}
  {45 minuti}
  {1 plumcake}
  {Facile}
  {Dolce da colazione}

\begin{ingredients}
  \ingredient{Uova}{3}
  \ingredient{Farina 00}{300 g}
  \ingredient{Zucchero}{100 g}
  \ingredient{Yogurt}{175 g}
  \ingredient{Olio di semi di girasole}{130 ml}
  \ingredient{Lievito vanigliato per dolci}{16 g}
  \ingredient{Farina di cocco}{q.b. per lo stampo}
  \ingredient{Olio di semi}{un goccio per lo stampo}
  \ingredient{Mirtilli, cioccolato o altri ingredienti a piacere}{opzionali}
\end{ingredients}

\begin{tools}
  \tool{Ciotola}
  \tool{Frusta}
  \tool{Stampo da plumcake}
  \tool{Carta assorbente}
\end{tools}

\begin{procedure}
  \step{Preriscalda il forno ventilato a 175\textdegree C.}

  \step{Versa un goccio di olio di semi sul fondo dello stampo da plumcake. Con carta assorbente distribuisci l'olio su tutto il fondo e sulle pareti dello stampo, poi cospargi con poca farina di cocco eliminando l'eccesso.}

  \step{Rompi 3 uova in una ciotola, aggiungi 100 g di zucchero e monta con la frusta per qualche minuto, fino a ottenere un composto piu' chiaro, leggermente gonfio e spumoso.}

  \step{Aggiungi 175 g di yogurt e 130 ml di olio di semi di girasole. Mescola con la frusta fino a ottenere una crema liscia ed emulsiona bene l'olio nell'impasto.}

  \step{Setaccia 300 g di farina 00 insieme a 16 g di lievito vanigliato per dolci. Aggiungili poco alla volta al composto di uova, zucchero, yogurt e olio, mescolando con la frusta solo il necessario per non sviluppare troppo glutine.}

  \step{Se vuoi arricchire il plumcake, incorpora delicatamente mirtilli, cioccolato o altri ingredienti a piacere. Mescola poco, solo per distribuirli nell'impasto.}

  \step{Versa l'impasto nello stampo da plumcake preparato, livella la superficie e inforna a 175\textdegree C in forno ventilato per circa 45 minuti.}

  \step{Controlla la cottura infilando uno stecchino al centro del plumcake: deve uscire asciutto o con poche briciole umide, ma senza impasto crudo.}

  \step{Lascia intiepidire il plumcake nello stampo per 10 minuti, poi sformalo e fallo raffreddare su una griglia prima di tagliarlo.}
\end{procedure}

\begin{notes}
  \note{Il plumcake e' super soffice ed e' ottimo per la colazione.}
  \note{Mirtilli, cioccolato, frutta secca o scorza di agrumi si possono aggiungere in base ai propri gusti.}
\end{notes}
